\chapter{Contexte général}

\section*{Introduction}
Ce chapitre introductif détaillera l'étude de l'existant, la problématique associée, ainsi qu'une évaluation approfondie de la solution proposée et des objectifs globaux du projet. Enfin, la dernière partie définira la méthodologie de gestion de projet adoptée.

\section{Présentation du cadre du projet}
Dans cette section, nous fournirons une analyse critique de l'existant, ainsi qu'une présentation détaillée de la solution proposée.
\subsection{Etude de l'existant}
Avant le développement de l'application Planora, une étude de l'existant a été réalisée afin d'analyser les solutions déjà disponibles sur le marché dans le domaine de la gestion de projets et du suivi collaboratif.Plusieurs plateformes et applications web proposent des fonctionnalités similaires, telles que la création de projets, la gestion des tâches, l'organisation des sprints, le suivi du backlog ou encore la collaboration entre les membres d'une équipe.\\
Parmi les solutions existantes, on retrouve des outils de gestion de projet comme Jira, Trello, Asana ou Monday.com, ainsi que des plateformes orientées vers les méthodes agiles et la communication d'équipe. Certaines de ces solutions intègrent également des tableaux de bord, des calendriers, des espaces de discussion, des systèmes de notification et des fonctionnalités d'analyse pour accompagner les équipes dans le suivi de leur avancement.\\
Cette étude a permis d'identifier les fonctionnalités communes, les bonnes pratiques ainsi que les limites techniques et fonctionnelles des solutions existantes, notamment en matière de centralisation des données, de suivi de la santé de l'équipe et d'intégration de fonctionnalités intelligentes d'aide à la décision.

\subsection{Tableau comparatif des solutions existantes}

Le tableau \ref{tab:comparaisonSolutions} présente une comparaison entre quelques solutions existantes de gestion de projet et la solution proposée Planora.

Dans ce tableau, la mention « Partiel » indique que la fonctionnalité existe de manière limitée ou nécessite une intégration avec un outil externe.

\begin{longtable}[c]{|p{0.14\textwidth}|p{0.22\textwidth}|p{0.20\textwidth}|p{0.20\textwidth}|p{0.18\textwidth}|}
  \caption{Comparaison entre les solutions existantes et Planora}
  \label{tab:comparaisonSolutions}                                                                                                                                                                                                                                                                                       \\

  \hline
  \textbf{Solution} & \textbf{Tâches, backlog et sprints}                                            & \textbf{Chat}                                       & \textbf{Réunions et calendrier}                              & \textbf{Analyse intelligente}                                                                \\
  \hline
  \endfirsthead

  \hline
  \textbf{Solution} & \textbf{Tâches, backlog et sprints}                                            & \textbf{Chat}                                       & \textbf{Réunions et calendrier}                              & \textbf{Analyse intelligente}                                                                \\
  \hline
  \endhead

  \hline
  \endlastfoot

  Jira              & Oui                                                                            & Partiel : via intégration avec Slack ou Teams       & Partiel : via intégration avec Google Calendar ou Outlook    & Non                                                                                          \\
  \hline
  Trello            & Partielle : gestion simple des cartes, sans sprints avancés natifs             & Partiel : via commentaires ou intégrations externes & Limitée : calendrier disponible selon les vues et extensions & Non                                                                                          \\
  \hline
  Asana             & Partielle : gestion des tâches avancée mais backlog et sprints non spécialisés & Partiel : via commentaires ou intégrations externes & Oui                                                          & Non                                                                                          \\
  \hline
  Monday.com        & Partielle : gestion personnalisable mais agile avancé selon configuration      & Oui                                                 & Oui                                                          & Non                                                                                          \\
  \hline
  Planora           & Oui                                                                            & Oui                                                 & Oui                                                          & Oui : analyse des messages pour détecter le stress, la frustration et l'ambiance de l'équipe \\
  \hline
\end{longtable}

\subsection{Critique de l'existant}
Malgré la diversité des solutions existantes, plusieurs insuffisances ont été constatées. La majorité des plateformes de gestion de projet proposent des fonctionnalités intéressantes, telles que la gestion des tâches, la planification des sprints ou le suivi de l'avancement, mais elles restent parfois dispersées entre plusieurs outils, ce qui complique le travail des équipes et limite la centralisation des informations.\\
De plus, certaines applications manquent de simplicité dans leur utilisation, notamment pour les équipes qui souhaitent adopter rapidement une méthode agile sans configuration complexe. L'absence d'un suivi précis de la santé de l'équipe constitue également une limite importante, car la plupart des solutions se concentrent principalement sur l'avancement technique du projet sans prendre en compte l'état général des membres, le stress ou la frustration exprimés pendant la collaboration.\\
Par ailleurs, les fonctionnalités de communication, de gestion des réunions et d'aide intelligente à la décision ne sont pas toujours intégrées dans une même plateforme, ce qui oblige les utilisateurs à utiliser plusieurs outils en parallèle. Ces limites soulignent la nécessité de concevoir une solution plus complète, centralisée et centrée sur les besoins réels des équipes de projet.

\subsection{Solution proposée}
Afin de remédier aux limites identifiées, le projet Planora propose une solution web complète regroupant plusieurs services de gestion de projet au sein d'une seule plateforme.\\
L'application permet aux utilisateurs de créer un compte sécurisé, de gérer des espaces de travail, de créer des projets, d'inviter des membres et d'organiser les tâches selon une approche collaborative et agile. Elle offre également la possibilité de gérer le backlog, les sprints, les sous-tâches, les réunions ainsi que le suivi de l'avancement à travers des tableaux de bord.\\
L'intégration d'un module basé sur l'intelligence artificielle constitue un élément différenciateur majeur, permettant d'analyser les messages échangés entre les membres afin d'évaluer la santé de l'équipe et de détecter les signes de stress ou de frustration. De plus, la plateforme intègre un espace de messagerie favorisant la communication et la coordination entre les membres du projet.\\
Cette solution vise à améliorer l'organisation, la collaboration et le suivi des projets tout en assurant une expérience utilisateur cohérente, centralisée et adaptée aux besoins réels des équipes.

\section{Choix Méthodologique}
Il est primordial d'opter pour la bonne méthodologie avant de débuter un projet informatique. En effet, aucune méthode universelle ne peut garantir à elle seule
l'exécution adéquate de l'ensemble des tâches d'un projet, tant les contextes et les objectifs varient d'un projet à l'autre.\\
La méthode sélectionnée doit donc être adaptée aux spécificités et aux contraintes de chaque projet. Ainsi, il est essentiel de retenir une approche qui réponde aux
besoins réels des parties prenantes, afin de garantir un résultat optimal et conforme aux attentes.\\
Dans le présent contexte, nous avons opté pour une approche de gestion de projet hybride, connue sous le nom de modèle \textbf{Water-Scrum-Fall}, qui combine les avantages des deux méthodologies les plus répandues :
\begin{itemize}
  \item La planification rigoureuse et la documentation exhaustive du \textbf{modèle
          en cascade} lors des phases d'audit et de conception, en s'appuyant sur \ac{UML}
        comme langage de modélisation.
  \item Le développement itératif et incrémental de la méthodologie \textbf{Scrum}
        lors des phases de mise en place et de développement, organisé en sprints permettant des livraisons régulières et une adaptation continue aux retours obtenus.
\end{itemize}

La figure \ref{MethHybride} ci-dessous illustre la méthodologie hybride adoptée :
\begin{figure}[H]
  \centering
  \includegraphics[width=0.6\columnwidth]{MethHybride.png}
  \caption{Méthodologie Hybride adoptée\cite{HybridMeth}}
  \label{MethHybride}
\end{figure}

En tirant parti des forces de chaque méthodologie, cette approche hybride vise à maximiser l'efficacité, à atténuer les risques et à assurer une livraison progressive et contrôlée des différents composants du projet.

\subsection{Méthodologie du modèle en cascade}
Le modèle en cascade, caractérisé par son approche séquentielle, est souvent associé lors de la phase de conception dans nombreux projets. Cette approche garantit une analyse approfondie des exigences et des objectifs du projet, suivie par la conception d'une solution complète et détaillée répondant à ces exigences. \\
Cette conception est ensuite examinée et validée avant de passer à la phase de développement proprement dite.\\
La linéarité du processus en cascade permet une gestion efficace des ressources et une planification minutieuse des étapes, offrant une structure et une prévisibilité précieuses dans la phase d'analyse des besoins et de conception\cite{Waterfall}
\subsection{Méthodologie agile}
La méthodologie Agile adoptée durant le cycle de vie de développement se distingue
des approches traditionnelles telles que le cycle en V ou le modèle en cascade, qui nécessitent une planification complète du projet avant le début du développement, par sa flexibilité et son adaptabilité en se concentrent sur des objectifs à court terme. \\
Plaçant le client au cœur du processus, cette approche favorise une collaboration étroite et des retours réguliers, ce qui permet à l'équipe de développement d'intégrer rapidement les modifications nécessaires pour répondre aux évolutions du projet\cite{MethAgile}.
La figure \ref{CycleDevAgile} ci-dessous illustre le cycle de vie de développement d'un logiciel agile :
\begin{figure}[H]
  \centering
  \includegraphics[width=0.85\columnwidth]{Agile Software Development Life Cycle.png}
  \caption{Cycle de vie de développement d'un logiciel agile}
  \label{CycleDevAgile}
\end{figure}
\subsubsection{Méthodologie Scrum}
La méthodologie Scrum est de loin la méthode Agile la plus utilisée dans le monde, offrant un cadre organisé pour superviser la création d'applications complexes. Cette approche intègre de différentes pratiques centralisées sur la flexibilité et la collaboration avec le client afin de s'adapter aux exigences changeantes tout en reposant sur la maintenance d'un « Backlog », qui décrit les changements ou améliorations à apporter sous forme des « user stories ».
\subsubsection{Caractéristiques de Scrum}
Les livraisons fréquentes sont une caractéristique clé de Scrum. Pour atteindre cet objectif, Scrum emploie des itérations de développement, connues sous le nom de « sprints », qui durent généralement de 2 à 4 semaines. Durant ces sprints, les opérations de développement traite activement les problématiques métier en cours, abordant les dossiers d'analyse et de conception afin de réaliser les phases de développement, de test et de déploiement\cite{ScrumCarac}.
\subsubsection{Rôles dans Scrum}
La méthodologie Scrum repose sur trois rôles essentiels :
\begin{itemize}
  \item Product Owner : chargé de déterminer et de donner la priorité aux éléments du backlog du produit. Celle-ci informe l'équipe de développement des besoins et des exigences du client et prend les décisions concernant les fonctionnalités à fournir.
  \item Scrum Master : garantir que l'équipe comprend et suit les principes et les méthodes de Scrum. Celui-ci facilite la communication des membres de l'équipe et soutient l'équipe dans l'accomplissement de ses objectifs.
  \item Équipe de développement : sont chargés de transférer les éléments du backlog produit en fonctionnalités livrables, celles-ci collaborent pour atteindre les objectifs du sprint et sont autonomes dans leur approche de travail.
\end{itemize}
La figure \ref{ProcScrum} ci-dessous illustre le processus de la méthode Scrum :
\begin{figure}[H]
  \centering
  \includegraphics[width=1\columnwidth]{ProcScrum.png}
  \caption{Processus de la méthode Scrum}
  \label{ProcScrum}
\end{figure}
\section{Langage de Modélisation}
Le choix du langage de modélisation dans le présent projet a été porté sur \ac{UML}. Il est justifié par sa standardisation, sa capacité à clarifier et à documenter les spécifications du système, ainsi que par son utilité dans l'analyse, la conception, la maintenance et la validation du système logiciel.\\
En effet, \ac{UML} offre un ensemble de conventions visuelles pour représenter les systèmes et les processus à l'aide de diagrammes et de représentations graphiques. Cette approche simplifie la communication des concepts complexes, favorisant ainsi une collaboration fluide et une compréhension approfondie à chaque étape du développement et de la conception.
\subsection{Présentation de UML}
Le langage de modélisation \ac{UML} est l'un des langages de modélisation les plus couramment employés dans le domaine de l'ingénierie logicielle. \ac{UML} propose une méthode efficace pour représenter visuellement la structure, le comportement et les interactions d'un système logiciel en utilisant des diagrammes standardisés tels que les diagrammes de classes, de cas d'utilisation, de séquence et bien d'autres\cite{UML}.
\subsection{Diagrammes UML}
Le langage de modélisation \ac{UML} est défini par 14 types de diagrammes, distingués en deux catégories : "structure" et "comportement", afin de garantir une modélisation appropriée des caractéristiques de tout système.
Afin de garantir une modélisation complète et précise des fonctionnalités de la future application, nous avons choisi d'utiliser les diagrammes UML répertoriés dans le tableau \ref{tab:diagUmlTab} ci-dessous :

\begin{longtable}[c]{|p{0.25\textwidth}|p{0.65\textwidth}|}
  \caption{Types des diagrammes UML}
  \label{tab:diagUmlTab}                                                                                                                                                                                                                                                                                                                                                                                                                               \\

  \hline
  \textbf{Diagramme}             & \textbf{Description}                                                                                                                                                                                                                                                                                                                                                                                                \\
  \hline
  \endfirsthead
  \hline
  \textbf{Diagramme}             & \textbf{Description}                                                                                                                                                                                                                                                                                                                                                                                                \\
  \hline
  \endhead
  \hline
  \endlastfoot
  Diagramme de classe            & Le diagramme UML le plus largement utilisé, essentiel à toute approche orientée objet, est le diagramme de classe. Ce diagramme représente des entités spécifiques du système applicatif sous forme de classes. Chaque classe est composée d'attributs, définissant les caractéristiques des objets créés à partir de ces classes, ainsi que de méthodes, définissant les actions que ces objets peuvent effectuer. \\
  \hline
  Diagramme de cas d'utilisation & Le diagramme de cas d'utilisation représente une fonctionnalité spécifique dans un système applicatif et met en évidence les liens et les interactions avec ses différentes fonctionnalités à l'aide d'acteurs internes et externes.                                                                                                                                                                                \\
  \hline
  Diagramme de séquence          &
  Le diagramme de séquence illustre comment les objets d'un scénario spécifique interagissent les uns avec les autres dans un ordre chronologique.                                                                                                                                                                                                                                                                                                     \\
  \hline
\end{longtable}

\section*{Conclusion}
Dans ce premier chapitre, nous avons présenté le contexte général du projet Planora, en mettant en lumière l'étude de l'existant, la critique des solutions disponibles ainsi que la solution proposée. Nous avons également exposé le choix méthodologique adopté, basé sur une approche hybride Water-Scrum-Fall, ainsi que le langage de modélisation UML retenu pour la conception du système. Dans le chapitre suivant, nous approfondirons l'analyse et la spécification des besoins fonctionnels et non fonctionnels de l'application.